% Part: model-theory
% Chapter: interpolation
% Section: introduction

\documentclass[../../../include/open-logic-section]{subfiles}

\begin{document}

\olfileid[tr]{mod}{int}{int}

\olsection{Giriş}

Aradeğerleme teoremi şu sonuçtur: $\Entails !A\lif !B$ varsayılsın.
O zaman $\Entails !A\lif !C$ ve $\Entails !C\lif !B$ olacak bir
!!a{sentence}~$!C$ vardır. Üstelik $!C$ içindeki her !!{constant},
!!{function} ve !!{predicate} ($\eq$ dışında) hem $!A$ hem $!B$ içinde
geçer. !!{sentence}~$!C$ ifadesine $!A$ ile $!B$'nin bir
\emph{aradeğerleyeni} denir.

Aradeğerleme teoremi başlı başına ilginçtir; fakat başlıca önemi, bir
kuramda tanımlanabilirlik hakkındaki sonuçları ve iki tutarlı kuramı
birleştirmenin hangi koşullarda tutarlı bir kuram verdiğini ispatlamak
için kullanılabilmesidir. İlk sonuç Beth tanımlanabilirlik teoremi,
ikincisi Robinson ortak tutarlılık teoremi olarak bilinir.

\end{document}
